%% kapa: exam=1 round=31 year=2020 period=1 form=A subject=economics subject_ko=경제학원론 seq=24 number=64
%% source: Q-Net 감정평가사 기출문제 (cst00309) artlSeq=5207219 file=extracted/1008/5207219/A형/A형/2020년도 제31회 감정평가사 1차 시험_1교시_A형.pdf
%% conversion: kordoc 4.24.51 -> tools/gen1.py (auto); review=proofread-sonnet; flags=none
\begin{question}{64}
소비이론에 관한 설명으로 옳지 않은 것은?
\begin{choices}
  \item 생애주기가설에 따르면 장기적으로 평균소비성향이 일정하다.
  \item 항상소득가설에 따르면 단기적으로 소득 증가는 평균소비성향을 감소시킨다.
  \item 케인즈(M. Keynes)의 소비가설에서 이자율은 소비에 영향을 주지 않는다.
  \item 피셔(I. Fisher)의 기간 간 소비선택이론에 따르면 이자율은 소비에 영향을 준다.
  \item 임의보행(random walk)가설에 따르면 소비의 변화는 예측할 수 있다.
\end{choices}
\end{question}
\answer{5}
